\section{Arquitectura Final del Sistema}
\label{sec:arquitectura_final_sistema}

La arquitectura del sistema desarrollado implementa un ecosistema modular capaz de manejar múltiples modalidades de datos biomédicos en tiempo real. Esta sección presenta la arquitectura final del sistema, describiendo los patrones arquitectónicos fundamentales y las decisiones de diseño que facilitan la operación independiente de los tres módulos principales.

El sistema final está compuesto por tres módulos independientes pero interoperables: \textbf{\textit{Input Collector}} para captura de eventos de usuario, \textbf{\textit{EEG Processor}} para procesamiento de señales cerebrales, y \textbf{\textit{API Explorer}} como \textit{\gls{framework}} web unificado para exploración y control remoto. Cada módulo implementa patrones arquitectónicos específicos optimizados para sus requisitos particulares, mientras mantiene la capacidad de exposición de funcionalidades mediante \acrshort{api} que facilitan la integración futura.

\subsection{Visión General de la Arquitectura}
\label{subsec:vision_general_arquitectura}

La arquitectura global del sistema se fundamenta en principios de diseño modular que permiten que cada módulo opere de forma completamente independiente. Esta independencia no impide la coordinación cuando es necesaria, sino que la facilita mediante interfaces de acceso a datos y capacidades de control remoto.

\subsubsection{Principios Arquitectónicos Fundamentales}
\label{subsubsec:principios_fundamentales}

El desarrollo del sistema ha consolidado un conjunto de principios arquitectónicos que se aplican consistentemente a través de todos los módulos, proporcionando coherencia estructural y facilitando tanto el mantenimiento como la extensibilidad futura.

\textbf{Arquitectura Dirigida por Eventos}: La comunicación interna de cada módulo utiliza un sistema de eventos centralizado que actúa como columna vertebral, eliminando dependencias directas entre componentes y facilitando la coordinación de operaciones complejas. Este patrón ha demostrado ser especialmente efectivo para el manejo de datos de alta frecuencia y operaciones asíncronas.

\textbf{Patrón Singleton Thread-Safe}: Los servicios centrales de cada módulo implementan el patrón \textit{singleton} con \textit{\gls{double_checked_locking}} para garantizar una instancia única con operaciones seguras en entornos multihilo. La \gls{inicializacion_perezosa} optimiza el uso de recursos retrasando la creación de objetos hasta que sean realmente necesarios.

\textbf{Separación de Responsabilidades}: La arquitectura implementa clara separación entre capas de presentación, lógica de negocio, acceso a datos, y abstracción de \textit{hardware}. Esta separación facilita significativamente el mantenimiento y la extensibilidad, como se evidenció durante la refactorización arquitectónica realizada en la segunda iteración del desarrollo.

\subsubsection{Flujos de Datos del Sistema}
\label{subsubsec:flujos_datos_sistema}

Cada módulo del sistema puede operar completamente independiente, generando y procesando sus propios flujos de datos. Esta aproximación permite que los investigadores utilicen los módulos según sus necesidades específicas sin requerir la activación completa del sistema.

\textbf{Grabación Independiente con Marcado Temporal}: Cada módulo puede grabar independientemente en formatos \acrshort{csv} con \textit{\glspl{timestamp}} precisos, permitiendo análisis correlacional posterior sin la sobrecarga de sincronización en tiempo real. Esta aproximación simplifica la arquitectura del sistema mientras proporciona capacidades de sincronización \textit{offline}.

\textbf{Control Remoto via \acrshort{api}}: Los módulos exponen una \acrshort{api} \acrshort{rest} que permite control externo y acceso a datos desde otras aplicaciones. \textit{\textbf{API Explorer}} proporciona una interfaz unificada para estas capacidades.

\textbf{Streaming de Eventos}: El sistema de eventos interno de cada módulo puede transmitir información hacia \textit{\textbf{API Explorer}} mediante \textit{\glspl{websocket}}, facilitando monitoreo remoto y visualización coordinada en tiempo real.

\subsection{Componentes Principales del Sistema}
\label{subsec:componentes_principales_sistema}

\subsubsection{Módulo \textit{Input Collector}}
\label{subsubsec:modulo_input_collector}

\textit{\textbf{Input Collector}} constituye el componente del sistema especializado en captura, procesamiento y visualización de eventos de entrada desde múltiples dispositivos en tiempo real. Su desarrollo se centró en resolver los desafíos únicos asociados con la captura de eventos de usuario con latencia mínima y precisión temporal.

\textbf{Enfoque Arquitectónico}: El módulo implementa una \gls{arquitectura_hexagonal} que separa efectivamente la lógica de dominio de las interfaces externas. Esta decisión arquitectónica permitió desarrollar soporte unificado para dispositivos como teclado, ratón y controladores \acrshort{ps4}, mientras se mantenían múltiples interfaces de acceso sin duplicación de funcionalidad.

La arquitectura dirigida por eventos demostró ser particularmente efectiva para el controlador \acrshort{ps4}, donde se implementó un \textit{event loop} \textbf{\acrshort{sdl}} dedicado que resolvió completamente los problemas de rendimiento y compatibilidad que habían limitado la funcionalidad en interfaces \acrshort{cli}. Esta implementación elimina las dependencias de \textit{\gls{timers}} \textbf{\acrshort{qt}} y proporciona captura con latencia mínima independiente del contexto de ejecución.

\textbf{Capacidades y Aplicaciones}: El sistema de captura soporta múltiples modos de seguimiento que se adaptan a diferentes requisitos de aplicación. Los modos de \textit{tracking} de ratón incluyen \textit{delta} para movimiento en bruto, \textit{angle} para direcciones discretas, \textit{axis} para direcciones cardinales, y \textit{joystick} para representación analógica con velocidad por eje.

El objetivo principal del módulo es la generación de \textit{datasets} de alta calidad con marcado temporal preciso que permita la sincronización posterior con señales cerebrales para el desarrollo de sistemas \acrshort{bci}. Los \textit{\glspl{timestamp}} con precisión de milisegundo facilitan el análisis correlacional entre patrones de entrada y actividad neuronal, habilitando estudios sobre \gls{imaginacion_motora}, análisis de comportamientos de usuario durante tareas específicas, y desarrollo de interfaces cerebro-computadora híbridas que combinen intenciones detectadas mediante \acrshort{eeg} con acciones explícitas del usuario.

\textbf{Características técnicas}:
\begin{itemize}
    \item Arquitectura modular con servicios especializados
    \item Soporte múltiple: teclado, ratón, controladores \acrshort{ps4}
    \item Sistema dirigido por eventos con optimizaciones de rendimiento
    \item Interfaces: \acrshort{gui} \textbf{PyQt5}, \acrshort{cli}, \acrshort{api} web
\end{itemize}

La documentación técnica detallada del módulo se encuentra en el \textbf{Anexo~\ref{chap:anexo_input_collector}}.

\subsubsection{Módulo \textit{EEG Processor}}
\label{subsubsec:modulo_eeg_processor}

\textit{\textbf{EEG Processor}} constituye el núcleo especializado del sistema para procesamiento de \glspl{biosenial} cerebrales, implementando una \gls{arquitectura_orientada_servicios} que integra múltiples fabricantes de dispositivos comerciales con algoritmos de análisis neurofisiológico avanzados.

\textbf{Arquitectura Orientada a Servicios}: El módulo adopta una arquitectura orientada a servicios que separa efectivamente las responsabilidades entre colección de datos, procesamiento de señales, análisis avanzado, y presentación de resultados. La integración con \textbf{BrainFlow} como capa de abstracción de \textit{hardware} proporciona soporte unificado para más de 20 fabricantes de dispositivos, facilitando significativamente el desarrollo de aplicaciones que necesiten funcionar con múltiples tipos de \textit{hardware} sin modificaciones específicas.

\textbf{Características técnicas}:
\begin{itemize}
    \item Arquitectura orientada a servicios con detección automática de canales
    \item Integración \textbf{Brainflow} con soporte para múltiples fabricantes
    \item Procesamiento avanzado: \gls{ica_analysis} (\acrshort{ica}), \gls{wavelet}, \textbf{MNE-Python}, \gls{filtros_digitales}, \gls{downsampling}
    \item Visualización científica en tiempo real con múltiples modalidades
\end{itemize}

La documentación técnica detallada del módulo se encuentra en el \textbf{Anexo~\ref{chap:anexo_eeg_processor}}.

\subsubsection{Módulo \textit{API Explorer}}
\label{subsubsec:modulo_api_explorer}

\textit{\textbf{API Explorer}} constituye el \textit{\gls{framework}} web unificado del sistema, proporcionando capacidades de exploración dinámica y control remoto de las funcionalidades de los demás módulos. Su desarrollo se centró en crear una plataforma web que facilite tanto la exploración interactiva como la integración con otros sistemas.

\textbf{Framework Web Modular}: El módulo implementa una arquitectura web completa basada en \textbf{FastAPI} con soporte \textit{\gls{websocket}} para comunicación en tiempo real. El sistema de registro dinámico permite que módulos externos expongan sus funcionalidades mediante interfaces web generadas automáticamente, mientras que la arquitectura de \textit{\gls{worker}} \textit{pools} gestiona eficientemente el streaming de datos de alta frecuencia con \textit{\gls{throttling}} adaptativo.

La arquitectura \textit{frontend} modular implementa \textit{\gls{managers}} especializados que coordinan visualización de datos mediante \textbf{Chart.js}, gestión de temas oscuro/claro, y formularios dinámicos que se adaptan automáticamente a los tipos de parámetros definidos por cada módulo.

\textbf{Características técnicas}:
\begin{itemize}
    \item \textit{Framework} web modular con \textbf{FastAPI}
    \item Sistema de registro dinámico de módulos
    \item Visualización tiempo real con \textit{\gls{websocket}} y optimizaciones de rendimiento
    \item Interfaz web inspirada en la documentación automática de \textbf{FastAPI} con múltiples mejoras: modo claro/oscuro, mejores explicaciones y facilidades para parámetros mediante \textit{sliders}, selectores \textit{combobox}, y plantillas \acrshort{json} con verificador de correctitud. Incluye visualizador de gráficos, disponible solo cuando el módulo lo solicita durante el registro, que proporciona una tabla actualizada dinámicamente o modo gráficas. En el modo gráficas se puede elegir entre mostrar todas las señales en un gráfico unificado, con escala logarítmica en el eje Y para evitar distorsión por valores atípicos, o utilizar gráficos independientes. Ambas opciones permiten filtrado por tipo de canal y capacidad de ocultar canales específicos
\end{itemize}

La documentación técnica detallada del módulo se encuentra en el \textbf{Anexo~\ref{chap:anexo_api_explorer}}.

\subsection{Integración y Extensibilidad del Sistema}
\label{subsec:integracion_extensibilidad_sistema}

La arquitectura modular facilita tanto la integración como la extensión futura del sistema. Para aplicaciones que requieren correlación entre eventos de entrada y señales cerebrales, los módulos pueden coordinarse temporalmente mediante la captura simultánea con \textit{\glspl{timestamp}} sincronizados que permitan análisis posterior.

\textit{\textbf{API Explorer}} proporciona un punto de control unificado que permite la coordinación de operaciones entre módulos cuando es necesario, mientras que las capacidades de registro independientes facilitan análisis posterior con sincronización basada en \textit{\glspl{timestamp}} de alta precisión.

\subsection{Consideraciones de Rendimiento}
\label{subsec:rendimiento_sistema}

El sistema implementa múltiples optimizaciones para manejar datos de alta frecuencia en tiempo real. Estas incluyen \textit{\gls{batching}} de eventos para reducir el \textit{\gls{overhead}} de comunicación, \textit{\gls{worker}} \textit{pools} para distribución de procesamiento, \textit{\gls{buffer}} circular con \gls{descarte_adaptativo} para degradación elegante del rendimiento, y \gls{serializacion} optimizada con \textbf{orjson} que proporciona mejoras de rendimiento de 2-5x comparado con la biblioteca estándar.

La gestión de memoria implementa límites de \textit{\gls{buffer}} configurables con limpieza automática, mientras que el diseño \textit{\gls{thread-safe}} utiliza \textit{\gls{locks}} granulares, señales \textbf{\acrshort{qt}} para comunicación segura entre hilos, y aislamiento de errores que previene la propagación de fallos entre componentes.

\subsection{Capacidades Actuales del Sistema}
\label{subsec:capacidades_actuales_sistema}

El sistema desarrollado proporciona un módulo de captura de eventos de teclado, movimientos y clics de ratón, y controladores con \textit{\glspl{timestamp}} precisos y un módulo de conexión intuitiva a dispositivos \acrshort{bci} mediante \textbf{Brainflow} con capacidad de aplicar filtros en tiempo real y grabación en \acrshort{csv} con \textit{\glspl{timestamp}} y posibilidad de añadir \gls{marcadores}. Cada módulo graba independientemente en \acrshort{csv}, también \acrshort{json} para el módulo de entrada, permitiendo sincronización \textit{offline} posterior.

Esta aproximación presenta ventajas específicas: evita sobrecarga \acrshort{cpu} extra en tiempo real, previene fallos relacionados con sincronización entre módulos al mantener independencia de cada componente, y no representa pérdida significativa de valor ya que, aunque existiera sincronización en tiempo real, no estaría desarrollada la funcionalidad para servir ese \textit{\gls{stream}} sincronizado en tiempo real a aplicaciones externas.

El sistema constituye una herramienta útil que facilita la investigación y desarrollo en interfaces cerebro-computadora mediante herramientas especializadas e interfaces accesibles.

\subsection{Extensibilidad y Mantenimiento}
\label{subsec:extensibilidad_mantenimiento}

La arquitectura del sistema facilita extensiones futuras mediante la implementación de principios \acrshort{solid}, patrones de extensión bien definidos, y separación clara de responsabilidades. El sistema de registro dinámico en \textit{\textbf{API Explorer}} permite la adición de nuevos módulos sin modificación del código existente.